\clearpage
\fancyhead[L]{\small CAGNOLI. STERNVERZEICHNISS.}
\hypertarget{native-501-lower}{}\label{native-501-lower}
\begin{center}
\small Göttingische gelehrte Anzeigen. Stück 77. S. 761..763. 14. Mai 1808.\\[3pt]
\rule{.72\linewidth}{.3pt}
\end{center}

\emph{Catalogue de \textup{501} étoiles suivi des tables relatives d'Aberration et de Nutation, par Antoine Cagnoli.} Modena 1807. 280 Seiten im Quartformat.

Dieses Sternverzeichnis war bereits im X. Band der \emph{Memorie di Matematica e di Fisica della Societa Italiana} abgedruckt und erscheint hier in verbesserter und erweiterter Form. Jenem ersten Abdruck waren auch historische Mitteilungen über die Entstehung des Katalogs beigegeben, aus denen wir hier einiges herausgreifen. \textsc{Cagnoli} hatte schon im Jahr 1783 an seinem damaligen Wohnort Paris begonnen, ein Sternverzeichnis aus eigenen Beobachtungen zu erstellen. In den Jahren 1788\mbox{. . .}1792 vollendete er die Arbeit dann mit denselben Instrumenten in Verona. Sein Plan war, nördlich vom Äquator in jeder Zone von einem halben Grad Ausdehnung in Deklination wenigstens vier zuverlässige Bestimmungen von Sternen zu geben, die in Rektaszension ungefähr 6 Stunden auseinanderlagen. Das erklärt die zunächst auffallende Auswahl der Sterne. Es sind keineswegs die hellsten in jedem Sternbild ausgewählt, vielmehr finden sich an ihrer Stelle oft Sterne bis zur sechsten Größenklasse im Verzeichnis. \textsc{Cagnoli} konnte sich allerdings nicht ganz an diesen Plan halten, teils weil nicht überall geeignete Sterne zu finden waren, teils weil ihn gesundheitliche Probleme und andere Störungen daran hinderten. \textsc{Cagnoli} wollte für sein Verzeichnis nichts von anderen übernehmen. Um die absoluten Rektaszensionen zu erhalten, bestimmte er daher die von $\alpha$ im Fuhrmann durch 24 direkte Vergleiche mit der Sonne und erhielt so ein Ergebnis, das bis auf $1''$ mit der neuesten Angabe von \textsc{Maskelyne} übereinstimmte. Für die gesamte Arbeit verwendete er einen beweglichen Quadranten von drei Fuß Radius, dessen achromatisches Fernrohr zwei Zoll Öffnung hat, und ein achromatisches Meridianfernrohr von 3$\frac{1}{2}$ Fuß mit 28 Linien Öffnung, beide von \textsc{Megnié}, sowie eine vortreffliche Pendeluhr, die \textsc{Robins} in Paris verfer%
\SourceLeaf{502}%
tigt hat. Alle diese Instrumente befinden sich gegenwärtig an der Sternwarte in Mailand. Einige Sterne sind übrigens nach Beobachtungen von \textsc{Cesaris} in Mailand bestimmt und im Katalog durch ein besonderes Zeichen gekennzeichnet.

Von den 501 Sternen im Verzeichnis sind die meisten nördliche, nur 28 sind südliche. Ihre Positionen sind für 1800 angegeben. Zu jeder Rektaszension und Deklination stehen die Anzahl der Beobachtungen und die Differenz der am stärksten voneinander abweichenden Beobachtungen. Bei den meisten Sternen betragen die Unterschiede gegenüber den Bestimmungen von \textsc{Piazzi} nur wenige Sekunden. Es zeugt wirklich von besonderem Geschick und großer Sorgfalt, dass \textsc{Cagnoli} mit seinen für einen solchen Zweck doch nur mittelmäßigen Hilfsmitteln diese Genauigkeit erreichen konnte. Am Ende des Verzeichnisses sind sämtliche Sterne noch einmal nach ihren Abständen vom Nordpol geordnet. Das ist für alle sehr bequem, die diese Sterne zum Vergleich heranziehen wollen.

Besonders wertvoll sind die dem neuen Abdruck des Katalogs beigefügten speziellen Aberrations- und Nutationstafeln für die meisten dieser Sterne. Die Aberrationstafeln sind bereits nach der neuesten Bestimmung der Aberrationskonstanten durch \textsc{Delambre} ($20''25$) berechnet. Die Nutationstafeln beruhen dagegen noch auf der Annahme, dass die große und die kleine Halbachse der Nutationsellipse $9''0$ und $6''7$ betragen. Alle ihre Zahlen müssen daher um ein Neuntel oder ein Vierzehntel erhöht werden, wenn sie mit den neuesten Angaben von \textsc{Laplace} oder \textsc{von Zach} in Übereinstimmung gebracht werden sollen.
